% !TeX root = ../../tesi.tex

\section{Architettura del sistema}

\subsection{Visione d’insieme}

Illustro l’architettura di ChatCart con un diagramma: il \emph{widget} nel negozio (\emph{app
layer}), il \emph{server} che ospita l’agente (\emph{agent layer}), la \emph{dashboard} del
commerciante e i servizi esterni. Motivo la scelta di uno stile funzionale anziché orientato agli
oggetti, la decisione, su indicazione del \emph{tutor}, di non introdurre astrazioni prima che
fossero necessarie e i \emph{design pattern} che ho adottato quando lo sono diventate.

\subsection{Il contratto tra agente e \emph{widget}}
\label{sec:contratto}

Approfondisco la scelta su cui si regge l’interfaccia generativa: un unico elenco di strumenti,
condiviso tra agente e \emph{widget}, che distingue gli strumenti che leggono i dati del negozio da
quelli che mostrano un componente e chiudono la risposta. Spiego come da questo elenco derivo i tipi
condivisi tra le due parti, così che il \emph{widget} sappia sempre disegnare ciò che l’agente
sceglie.
